\documentclass[11pt,a4paper]{article}

% ---------- Packages ----------
\usepackage[margin=1in]{geometry}
\usepackage{times}
\usepackage[T1]{fontenc}
\usepackage[utf8]{inputenc}
\usepackage{amsmath,amssymb}
\usepackage{booktabs}
\usepackage{longtable}
\usepackage{array}
\usepackage{ragged2e}
\usepackage{enumitem}
\usepackage{hyperref}
\usepackage{titlesec}
\usepackage{fancyhdr}
\usepackage{abstract}

\hypersetup{
  colorlinks=true,
  linkcolor=black,
  citecolor=black,
  urlcolor=blue,
  pdftitle={AI in Finance: Stock Market Prediction Using AI and ML},
  pdfauthor={Sanjana Singh}
}

\titleformat{\section}{\large\bfseries}{\thesection.}{0.6em}{}
\titleformat{\subsection}{\normalsize\bfseries}{\thesubsection}{0.6em}{}

\pagestyle{fancy}
\fancyhf{}
\fancyhead[C]{\small AI in Finance: Stock Market Prediction Using AI and ML}
\fancyfoot[C]{\thepage}
\renewcommand{\headrulewidth}{0.2pt}

\setlength{\parindent}{0pt}
\setlength{\parskip}{0.6em}

% Custom column type for wrapped table cells
\newcolumntype{L}[1]{>{\RaggedRight\arraybackslash}p{#1}}

\begin{document}

% ---------- Title ----------
\begin{center}
  {\LARGE\bfseries AI in Finance: Stock Market Prediction Using AI and ML}\\[0.4em]
  {\large A Systematic Literature Review}\\[1em]
  {\normalsize Sanjana Singh}\\[0.2em]
  {\small End-Sem Research Methodology Project [SEM: 6]}
\end{center}

\vspace{1em}

% ---------- Abstract ----------
\begin{abstract}
\noindent This paper presents a systematic review of the literature on the application of machine
learning and artificial intelligence to predict stock market movements. The financial market process
is highly volatile and non-linear. Therefore, several studies have attempted to use advanced data
science techniques such as Random Forest, Support Vector Machine, Artificial Neural Network, Long
Short-Term Memory, and hybrid models to predict stock market trends with better accuracy. The
results indicate that while the use of artificial intelligence and machine learning leads to better
predictions, higher accuracy is still challenging mainly because of the volatility of the market. The
study recommends hybrid and ensemble models to better analyze multiple variables influencing stock
market trends. This review highlights research opportunities and directions for future studies on
the topic.
\end{abstract}

% =====================================================
\section{Introduction}

Right from the start, investor mood shapes how markets move, showing what people think about
future profits \cite{khan2022}. Instead of steady patterns, price shifts often jump around without
clear cause, reacting fast to news or rumors \cite{subasi2021}. Money flows into stocks because
companies need funds to grow, which helps entire economies expand over time. Still, small changes
in data or policy can ripple through trading activity in ways that surprise even experts. Over the
years, scholars have looked closely at trends using mathematical tools, trying to spot hidden order
beneath chaos. Globally, institutions watch these swings carefully since outcomes affect jobs, spending,
and savings alike. Even with powerful computers, predictions fail regularly due to too many moving
parts influencing results. Because emotions mix with facts during trades, logic alone rarely explains
sudden market turns \cite{rouf2021}.

Back then, picking stocks meant digging into numbers or studying charts. Looking at balance sheets,
profits, and income streams -- this way tries to find what a business is truly worth. Charts tell
another story, one built on past prices and how much was traded, searching for familiar rhythms
\cite{subasi2021}. Even though these methods work sometimes, today's markets spit out too much
information too fast for them to keep up. When emotions swing wildly or something unexpected hits,
those old tools often fall short. Their blind spots grow bigger when change arrives without warning
\cite{kumar2022}.

Twenty years of fast progress in artificial intelligence together with machine learning have reshaped
how forecasts are made in the stock market. Traditional number-crunching approaches fall short where
these systems thrive -- spotting hidden links inside massive piles of information \cite{rouf2021}. They
learn from what happened before, adjusting when markets shift direction. Instead of relying on one
method, researchers mix different tools like neural nets, support vector machines, random forest
setups, LSTM cells, or combined strategies. Each handles specific kinds of math challenges tied to
pricing trends and risk signals. Deeper layers of processing now pull meaning from messy sources such
as tweets or headlines through emotion tracking tricks baked into advanced networks \cite{khan2022}.
These upgrades let predictions run deeper than older formulas ever allowed.

Despite the proliferation of research in this field, there is no consensus on a single ``optimal'' method
for stock market prediction. Performance varies significantly depending on the choice of algorithm,
data preprocessing, feature selection, and evaluation metrics \cite{ferreira2021}. Moreover, while
numerous studies introduce specific models, comprehensive reviews that synthesize findings across
different methodologies and identify future research directions remain relatively scarce
\cite{sizan2023,soni2022}.

This paper provides a systematic review of recent research on AI and machine learning in stock market
forecasting. It categorizes the various methodologies employed, examines the types of data used, and
analyzes the performance metrics reported in the literature. By identifying the strengths and
weaknesses of current approaches, this work aims to highlight critical research gaps and suggest
directions for future study \cite{rouf2021,sizan2023}.

The remainder of this paper is organized as follows. Section 2 presents the background and
foundational concepts relevant to the topic. Section 3 reviews the machine learning methodologies
identified in the literature. Section 4 provides a comparative analysis of the selected studies. Section
5 identifies key research gaps. Section 6 suggests directions for future work. Section 7 concludes the
paper.

% =====================================================
\section{Background}

The stock market is a dynamic and essential component of the global financial system, reflecting
economic health and providing opportunities for wealth accumulation \cite{strader2020}. However,
predicting price movements remains a formidable challenge due to the high degree of uncertainty and
the influence of diverse factors, including corporate earnings, geopolitical events, and investor
psychology \cite{khan2022,subasi2021}. The complexity of these interactions makes the market
inherently difficult to forecast \cite{kumar2022}.

Traditional forecasting methods are generally categorized into fundamental and technical analysis.
Fundamental analysis focuses on economic and financial indicators to determine a stock's intrinsic
value \cite{mokhtari2021}. Technical analysis, on the other hand, utilizes historical price patterns
and indicators such as the Relative Strength Index (RSI), Moving Average Convergence Divergence
(MACD), and Simple Moving Averages (SMA) to predict future price movements \cite{subasi2021}.
While fundamental analysis seeks to determine what a stock should be worth, technical analysis
attempts to identify where the price is likely to go based on past trends.

Two major theories have historically guided the understanding of market predictability: the Efficient
Market Hypothesis (EMH) and the Adaptive Market Hypothesis (AMH). EMH suggests that stock
prices fully reflect all available information, making it impossible to consistently outperform the
market \cite{chopra2021}. In contrast, AMH posits that market efficiency is not static but evolves
over time, allowing for periods of predictability when certain conditions are met \cite{mokhtari2021}.

The availability of high-performance computing and massive financial datasets has shifted the focus
toward AI-driven forecasting \cite{subasi2021}. These systems can identify patterns in historical data
without the need for explicit programming, making them well-suited for the complex and non-linear
nature of financial markets \cite{sonkavde2023}. Machine learning tasks in this domain typically
involve either classification (e.g., predicting ``buy,'' ``sell,'' or ``hold'' signals) or regression (e.g.,
predicting specific price values) \cite{akhtar2022}.

Prominent techniques in machine learning research include Artificial Neural Networks (ANN),
Support Vector Machines (SVM), and Random Forests. Additionally, Long Short-Term Memory
(LSTM) networks -- a specialized type of Recurrent Neural Network (RNN) -- have become popular
for their ability to capture long-term dependencies in time-series data \cite{sheth2023,lin2024}.
Hybrid and ensemble methods, which combine the outputs of multiple models, are also increasingly
used to improve forecast reliability and reduce error \cite{sonkavde2023}.

Recent research has also begun to incorporate sentiment analysis, using NLP to extract market mood
from social media and financial news. This multi-source approach represents a significant shift in
stock market forecasting, as it acknowledges the impact of public opinion on trading behavior
\cite{khan2022,koukaras2022}.

% =====================================================
\section{Literature Review}

Over the past two decades, there has been a significant increase in the application of artificial
intelligence for stock market forecasting. This section reviews selected studies categorized by their
primary methodologies: classical machine learning, deep learning, hybrid/ensemble methods, and
multi-source sentiment analysis.

\subsection{Classical Machine Learning Methods}

Early research in AI-driven forecasting established the baseline for modern techniques by testing
fundamental machine learning models. Mokhtari et al. (2021) conducted a broad assessment using
technical indicators and financial fundamentals to predict Apple Inc. (AAPL) stock prices. Their
study compared several models, including Logistic Regression, K-Nearest Neighbors (KNN), Random
Forest, and SVM. Notably, Linear Regression achieved high accuracy for price prediction, with $R^2$
values nearing ideal levels, outperforming LSTM models in that specific context due to lower error
metrics (MAPE of 1.56 vs. 2.99). However, for sentiment classification, SVM proved most effective,
reaching an accuracy of 75.5\% and an AUC of 0.76. Despite these results, the authors noted that AI
models still face challenges in maintaining consistent performance over extended periods
\cite{mokhtari2021}.

Akhtar et al. (2022) examined the performance of Random Forest and SVM on historical price data,
reporting an aggregate accuracy of 80.3\%. Their findings emphasized that Random Forest slightly
outperformed SVM (80.8\% vs. 78.7\%) and highlighted the critical importance of data preprocessing.
Without rigorous cleaning and normalization, model stability and forecast accuracy were significantly
diminished \cite{akhtar2022}.

Sheth and Shah (2023) provided a detailed comparison of ANN, SVM, and LSTM. Their results showed
that ANNs consistently achieved approximately 90\% accuracy due to their ability to detect complex
non-linear trends. SVM performance varied between 60\% and 91\%, depending on the kernel selection
and parameter tuning. Interestingly, integrating SVM with a Genetic Algorithm for feature selection
improved precision to 93.7\%, demonstrating the value of optimization techniques \cite{sheth2023}.

\subsection{Deep Learning Models}

As deep learning evolved, LSTM networks became a preferred choice for financial time-series
forecasting. LSTMs address the vanishing gradient problem inherent in traditional RNNs, allowing
them to capture long-term dependencies in sequential data \cite{sonkavde2023}.

Sonkavde et al. (2023) analyzed various machine learning and deep learning models using Indian
stock datasets (TAINIWALCHM and AGROPHOS). While LSTM showed promise, XGBoost achieved
superior results in their sample, with an $R^2$ of 0.984 and a lower RMSE. The study concluded that
while LSTMs are powerful for temporal patterns, their performance is highly dependent on large
volumes of data and precise hyperparameter tuning \cite{sonkavde2023}.

Sheth and Shah (2023) further discussed the architectural advantages of LSTMs, noting that their
gated cell structure -- comprising input, forget, and output gates -- allows for selective information
retention. Their review reported LSTM accuracies ranging from 57\% to over 90\%, with
Attention-based LSTM variants often producing the most reliable trading signals \cite{sheth2023}.

\subsection{Ensemble and Hybrid Methods}

To overcome the limitations of individual algorithms, researchers have increasingly developed
ensemble and hybrid models. Sonkavde et al. (2023) proposed an ensemble framework combining
Random Forest, XGBoost, and LSTM. This hybrid model outperformed all individual algorithms,
achieving an $R^2$ of 0.992. The authors attributed this success to the model's ability to
simultaneously handle high-dimensional data, sequential learning, and temporal dependencies
\cite{sonkavde2023}.

Venkatarathnam et al. (2024) explored hybrid modeling using Genetic Algorithms (GA) to optimize
SVM and ANN parameters for NSE-listed stocks. Their results confirmed that GA-optimized models
consistently outperformed standalone versions. Additionally, they applied the ARIMA model for
short-term forecasting, finding that while it was effective for immediate trends, its accuracy declined
significantly over longer time horizons \cite{venkatarathnam2024}.

\subsection{Sentiment Analysis and Multi-Source Approaches}

A growing body of research incorporates textual and sentiment-based information, recognizing the
influence of news and public opinion on market movements. Mokhtari et al. (2021) demonstrated that
Twitter sentiment could serve as a proxy for fundamental analysis, though they noted the difficulty of
converting unstructured text into precise trading signals \cite{mokhtari2021}.

Sonkavde et al. (2023) highlighted that combining sentiment data with historical prices consistently
improves model accuracy. Training LSTMs on both price values and sentiment scores (encoded via
Word2Vec or N-Grams) led to substantial performance gains. They also emphasized the necessity of
filtering ``fake news'' and spam to maintain data quality \cite{sonkavde2023}.

\subsection{Meta-Review Findings}

Lin and Marques (2024) synthesized findings from ten major literature reviews, covering over 379
individual studies. Their meta-analysis confirmed that SVM, LSTM, and ANN remain the most widely
used AI methods in the field. Historical closing prices are the dominant input source (used in 66\% of
studies), followed by technical indicators and sentiment data. A key gap identified was the lack of
studies that simultaneously integrate social media sentiment, financial news, historical prices, and
macroeconomic data, suggesting that such multi-source hybrid approaches represent the next frontier
in research \cite{lin2024}.

% =====================================================
\section{Comparative Analysis}

Table~\ref{tab:comparative} summarizes the technique, dataset, reported accuracy, contribution, and
limitation of sixteen selected studies reviewed in this paper.

{\footnotesize
\setlength{\tabcolsep}{3pt}
\begin{longtable}{L{2.0cm} L{2.2cm} L{2.4cm} L{2.6cm} L{2.8cm} L{2.5cm}}
\caption{Comparative analysis of selected studies on AI/ML-based stock market prediction.}
\label{tab:comparative} \\
\toprule
\textbf{Paper} & \textbf{Technique} & \textbf{Dataset} & \textbf{Accuracy} & \textbf{Contribution} & \textbf{Limitation} \\
\midrule
\endfirsthead
\multicolumn{6}{l}{\textit{(continued from previous page)}} \\
\toprule
\textbf{Paper} & \textbf{Technique} & \textbf{Dataset} & \textbf{Accuracy} & \textbf{Contribution} & \textbf{Limitation} \\
\midrule
\endhead
\midrule
\multicolumn{6}{r}{\textit{continued on next page}} \\
\endfoot
\bottomrule
\endlastfoot

Sizan et al. (2023) \cite{sizan2023} & Random Forest, Gradient Boosting, Logistic Regression &
Historical stock prices (Yahoo Finance, Google Finance), financial indicators, news data & Not
specified & Provides a structured comparison of ML models for U.S. stock forecasting and highlights
practical implications for investors & Limited exploration of deep learning models; results depend
heavily on selected features \\

Sonkavde et al. (2023) \cite{sonkavde2023} & ML \& DL (Random Forest, XGBoost, LSTM, ensemble
models) & Real stock datasets (TAINIWALCHM, AGROPHOS) + financial features & Ensemble (RF +
XGBoost + LSTM) outperformed individual models & Proposes hybrid ensemble framework and
systematic review of ML/DL models & High computational complexity; interpretability is low \\

Akhtar et al. (2022) \cite{akhtar2022} & Random Forest, Support Vector Machine (SVM) & Historical
stock price dataset (pre-processed) & $\sim$80.3\% prediction accuracy & Emphasizes importance of
preprocessing and model tuning for better performance & Limited generalization due to dataset size
and lack of external factors \\

Rouf et al. (2021) \cite{rouf2021} & ML, text analytics, ensemble methods & Literature datasets
(2011--2021 studies) & Not applicable (survey-based) & Provides decade-long overview and identifies
emerging trends like sentiment analysis & No experimental validation; lacks real-world performance
metrics \\

Kumar \& Sarangi (2022) \cite{kumar2022} & ANN, SVM, Neural Networks & Review of 30 research
studies & Not specified & Identifies ANN and NN as most widely used techniques & No quantitative
comparison; lacks experimental benchmarking \\

Bansal et al. (2022) \cite{bansal2022} & KNN, Linear Regression, SVR, Decision Tree, LSTM & Indian
stock market data (12 companies, 7 years) & LSTM produced highest accuracy among models &
Provides empirical comparison between ML and DL methods on real-world data & Dataset limited
geographically; may not generalize globally \\

Khan et al. (2022) \cite{khan2022} & ML + Deep Learning + Sentiment Analysis & Social media
(Twitter), financial news, stock data & Up to 83.22\% (Random Forest ensemble) & Demonstrates
importance of sentiment data in improving prediction accuracy & Data noise and bias from social
media; requires heavy preprocessing \\

Sheth \& Shah (2023) \cite{sheth2023} & ANN, SVM, LSTM & General stock datasets & ANN showed
best results for nonlinear patterns & Explains strengths of neural networks in handling complex
relationships & LSTM requires large datasets; computational cost is high \\

Lin \& Marques (2024) \cite{lin2024} & SVM, LSTM, ANN (systematic review of reviews) & 379+
studies & Not applicable & Identifies dominant AI techniques and highlights research gaps & No
practical experimentation \\

Strader et al. (2020) \cite{strader2020} & ANN, SVM, Genetic Algorithms, Hybrid models &
Literature-based & Not applicable & Categorizes ML techniques into major research groups and
proposes future directions & Focuses on classification rather than performance comparison \\

Chopra \& Sharma (2021) \cite{chopra2021} & Neural networks, hybrid AI techniques & 148 research
papers & Not applicable & Provides taxonomy of AI techniques and research agenda & No empirical
validation; theoretical focus \\

Subasi et al. (2021) \cite{subasi2021} & Multiple ML classifiers & NASDAQ, NYSE, FTSE, Nikkei
datasets & Varies across models & Compares ML models across global markets & Performance
sensitive to dataset variations \\

Kompella \& Chilukuri (2020) \cite{kompella2020} & Random Forest + Sentiment Analysis &
Historical stock data + news sentiment & Not specified & Shows effectiveness of Random Forest
combined with sentiment data & Cannot achieve consistent accuracy due to volatility \\

Ampomah et al. (2021) \cite{ampomah2021} & Gaussian Na\"ive Bayes + PCA/LDA & Stock price
datasets & GNB + LDA performed best & Explores feature scaling and extraction impact on ML
performance & Limited comparison with advanced ML/DL models \\

Koukaras et al. (2022) \cite{koukaras2022} & ML + Sentiment Analysis (SVM, KNN, NB, etc.) &
Twitter, StockTwits, Yahoo Finance & F-score: 76.3\%, AUC: 67\% & Integrates sentiment analysis with
ML for better predictions & Dependent on sentiment data quality \\

Venkatarathnam et al. (2024) \cite{venkatarathnam2024} & AI \& ML techniques & General financial
datasets & Not specified & Highlights real-world impact of AI in stock trading & Lacks experimental
results \\

\end{longtable}
}

% =====================================================
\section{Research Gaps}

Despite significant advancements in AI-driven stock market prediction, several critical research gaps
remain. Addressing these limitations is essential for improving model robustness, generalizability,
and practical utility in real-world trading environments \cite{rouf2021,chopra2021}.

\subsection*{1. Limited Multi-Source Data Integration}
A significant portion of current research relies heavily on historical price data and a narrow set of
technical indicators. While these features are readily available, they represent only a fraction of the
factors influencing market dynamics. Fundamental data -- such as earnings reports, balance sheets,
and macroeconomic indicators -- is frequently neglected due to its unstructured and inconsistent
nature. Furthermore, although sentiment analysis of social media and news has gained traction, the
simultaneous integration of technical, fundamental, and sentiment data within a single predictive
framework remains rare. Developing models that can effectively synthesize these diverse data types is
a key area for future exploration \cite{sonkavde2023,lin2024}.

\subsection*{2. Overfitting and Lack of Generalizability}
A recurring challenge in the literature is the tendency of complex models, such as LSTMs and deep
ensembles, to overfit historical patterns, leading to poor performance when market conditions shift
\cite{sheth2023,bansal2022}. Many studies validate their models on limited datasets or specific time
periods without testing for cross-market or cross-temporal generalizability. Consequently, the
predictive power of these models in unseen market regimes remains largely unproven
\cite{subasi2021}.

\subsection*{3. Absence of Standardized Evaluation Metrics}
The lack of a unified evaluation framework complicates the comparison of results across different
studies. Researchers employ a wide array of metrics, including accuracy, F1-score, RMSE, MAE, and
$R^2$, often focusing on different assets or indices. This inconsistency makes it difficult to establish
benchmarks and verify the relative superiority of emerging techniques \cite{strader2020,lin2024}.

\subsection*{4. Short-Term vs. Long-Term Forecasting}
Most existing research focuses on short-term price movements, such as daily or weekly forecasts.
While valuable for high-frequency trading, these models are less applicable to long-term investment
planning. Developing reliable AI methodologies for extended time horizons remains a significant
challenge due to the increasing noise and uncertainty associated with long-term market trends
\cite{sonkavde2023,bansal2022}.

\subsection*{5. Under-Researched Emerging Markets}
The majority of literature focuses on developed markets, particularly in the U.S., China, and Western
Europe. Emerging economies in South Asia, Africa, and Latin America are often overlooked, despite
their growing influence in the global financial system. These markets present unique challenges,
such as lower liquidity and less consistent data reporting, which require tailored AI approaches
\cite{strader2020,venkatarathnam2024}.

\subsection*{6. Incomplete Exploration of Hybrid and Ensemble Methods}
While hybrid models often outperform standalone algorithms, many potential combinations remain
unexplored. For instance, integrating statistical models like ARIMA with deep learning architectures
like LSTM, or utilizing Genetic Algorithms for hyperparameter optimization in complex neural
networks, offers promising avenues for improving forecast accuracy that have yet to be fully realized
\cite{rouf2021,sonkavde2023}.

\subsection*{7. Neglect of Transaction Costs and Market Constraints}
Many models report high theoretical accuracy but fail to account for practical trading constraints,
such as transaction fees, slippage, and liquidity. In a real-world setting, these factors can significantly
erode predicted profits. Future research must incorporate realistic trading simulations to validate the
economic viability of AI-driven strategies \cite{ferreira2021}.

\subsection*{8. Model Interpretability and the ``Black Box'' Problem}
Advanced deep learning models often function as ``black boxes,'' providing little insight into the
underlying reasoning for their predictions. In the financial sector, where risk management and
regulatory compliance are paramount, the lack of interpretability is a major barrier to adoption.
Developing ``Explainable AI'' (XAI) techniques tailored for financial forecasting is crucial for
building trust among investors and regulators \cite{ferreira2021,chopra2021}.

% =====================================================
\section{Future Directions}

Future research should focus on enhancing the practical applicability and robustness of AI models in
financial markets. Key directions include:

\begin{itemize}[leftmargin=1.4em]
  \item \textbf{Integration of Alternative Data Sources:} Beyond traditional price and volume data,
  incorporating alternative sources such as satellite imagery, web search trends, and corporate
  filings could provide a more comprehensive view of market drivers.
  \item \textbf{Long-Term Forecasting Models:} Shifting focus toward multi-month and multi-year
  horizons will be beneficial for institutional investors and long-term financial planning.
  \item \textbf{Cross-Market Analysis:} Expanding research to include emerging markets will help
  validate the universality of AI models and uncover region-specific market dynamics.
  \item \textbf{Explainable AI (XAI):} Developing models that offer transparent reasoning for their
  predictions is essential for regulatory compliance and investor confidence.
  \item \textbf{Realistic Backtesting:} Incorporating transaction costs, market impact, and liquidity
  constraints into model evaluations will bridge the gap between theoretical research and actual
  trading performance.
\end{itemize}

The continuous evolution of AI and machine learning presents significant opportunities for the
financial sector. By addressing the identified research gaps and focusing on these future directions,
researchers can develop more reliable and transparent tools for stock market prediction
\cite{sonkavde2023,lin2024}.

% =====================================================
\section{Conclusion}

This systematic review has explored the application of artificial intelligence and machine learning in
stock market prediction, synthesizing insights from numerous academic studies. The literature
reveals a clear transition from traditional statistical methods to sophisticated deep learning and
ensemble architectures. While the inherent volatility and complexity of financial markets remain
significant challenges, the integration of advanced algorithms has demonstrated superior performance
in identifying non-linear patterns within large datasets \cite{ferreira2021,bansal2022}.

Key findings indicate that ensemble models -- particularly those combining Random Forest, XGBoost,
and LSTM -- consistently achieve higher accuracy by leveraging the complementary strengths of
different techniques \cite{sizan2023,sonkavde2023}. Furthermore, the inclusion of sentiment analysis
from news and social media has been shown to enhance predictive power, suggesting that market
sentiment is a critical factor alongside historical price data \cite{khan2022,koukaras2022}.

Despite these advancements, several hurdles persist. Many models suffer from overfitting and lack
generalizability across different market regimes. The absence of standardized evaluation metrics and
the ``black box'' nature of deep learning architectures further complicate the practical adoption of AI
in finance. Future research must address these issues by focusing on multi-source data integration,
long-term forecasting, and the development of explainable AI frameworks.

In conclusion, while no single method provides a definitive solution for stock market forecasting, the
continued development of AI and machine learning offers promising pathways for improving market
analysis and investment strategies. By addressing current research gaps and emphasizing model
transparency and robustness, the field can move toward more reliable and actionable financial
predictions \cite{strader2020,chopra2021,lin2024}.

% =====================================================
\begin{thebibliography}{99}

\bibitem{khan2022} Wasiat Khan, Mustansar Ali Ghazanfar, Muhammad Awais Azam, Amin Karami,
Khaled H. Alyoubi, and Ahmed S. Alfakeeh. Stock market prediction using machine learning
classifiers and social media, news. \textit{Journal of Ambient Intelligence and Humanized
Computing}, 13(7):3433--3456, 2022.

\bibitem{subasi2021} Abdulhamit Subasi, Faria Amir, Kholoud Bagedo, Asmaa Shams, and Akila
Sarirete. Stock market prediction using machine learning. \textit{Procedia Computer Science},
194:173--179, 2021.

\bibitem{rouf2021} Nusrat Rouf, Majid Bashir Malik, Tasleem Arif, Sparsh Sharma, Saurabh Singh,
Satyabrata Aich, and Hee-Cheol Kim. Stock market prediction using machine learning techniques: a
decade survey on methodologies, recent developments, and future directions. \textit{Electronics},
10(21):2717, 2021.

\bibitem{kumar2022} Deepak Kumar, Pradeepta Kumar Sarangi, and Rajit Verma. A systematic review
of stock market prediction using machine learning and statistical techniques. \textit{Materials
Today: Proceedings}, 49:3187--3191, 2022.

\bibitem{ferreira2021} Fernando G. D. C. Ferreira, Amir H. Gandomi, and Rodrigo T. N. Cardoso.
Artificial intelligence applied to stock market trading: a review. \textit{IEEE Access}, 9:30898--30917,
2021.

\bibitem{sizan2023} Mir Mohtasam Hossain Sizan, Bimol Chandra Das, Reza E. Rabbi Shawon, MD
Sohel Rana, Md Abdullah Al Montaser, Anchala Chouksey, and Laxmi Pant. AI-enhanced stock market
prediction: evaluating machine learning models for financial forecasting in the USA. \textit{Journal of
Business and Management Studies}, 5(4):152--166, 2023.

\bibitem{soni2022} Payal Soni, Yogya Tewari, and Deepa Krishnan. Machine learning approaches in
stock price prediction: a systematic review. In \textit{Journal of Physics: Conference Series}, volume
2161, page 012065. IOP Publishing, 2022.

\bibitem{strader2020} Troy J. Strader, John J. Rozycki, Thomas H. Root, and Yu-Hsiang John Huang.
Machine learning stock market prediction studies: review and research directions. \textit{Journal of
International Technology and Information Management}, 28(4):63--83, 2020.

\bibitem{mokhtari2021} Sohrab Mokhtari, Kang K. Yen, and Jin Liu. Effectiveness of artificial
intelligence in stock market prediction based on machine learning. \textit{arXiv preprint
arXiv:2107.01031}, 2021.

\bibitem{chopra2021} Ritika Chopra and Gagan Deep Sharma. Application of artificial intelligence in
stock market forecasting: a critique, review, and research agenda. \textit{Journal of Risk and
Financial Management}, 14(11):526, 2021.

\bibitem{sonkavde2023} Gaurang Sonkavde, Deepak Sudhakar Dharrao, Anupkumar M. Bongale,
Sarika T. Deokate, Deepak Doreswamy, and Subraya Krishna Bhat. Forecasting stock market prices
using machine learning and deep learning models: a systematic review, performance analysis and
discussion of implications. \textit{International Journal of Financial Studies}, 11(3):94, 2023.

\bibitem{akhtar2022} Md Mobin Akhtar, Abu Sarwar Zamani, Shakir Khan, Abdallah Saleh Ali Shatat,
Sara Dilshad, and Faizan Samdani. Stock market prediction based on statistical data using machine
learning algorithms. \textit{Journal of King Saud University -- Science}, 34(4):101940, 2022.

\bibitem{sheth2023} Dhruhi Sheth and Manan Shah. Predicting stock market using machine learning:
best and accurate way to know future stock prices. \textit{International Journal of System Assurance
Engineering and Management}, 14(1):1--18, 2023.

\bibitem{lin2024} Chin Yang Lin and Jo\~ao Alexandre Lobo Marques. Stock market prediction using
artificial intelligence: a systematic review of systematic reviews. \textit{Social Sciences \& Humanities
Open}, 9:100864, 2024.

\bibitem{koukaras2022} Paraskevas Koukaras, Christina Nousi, and Christos Tjortjis. Stock market
prediction using microblogging sentiment analysis and machine learning. In \textit{Telecom}, volume
3, pages 358--378. MDPI, 2022.

\bibitem{venkatarathnam2024} N. Venkatarathnam, Laxmana Rao Goranta, P. C. Kiran, B. P. G. Raju,
Samadhi Dilli, S. Mahabub Basha, and Manyam Kethan. An empirical study on implementation of AI
\& ML in stock market prediction. \textit{Indian Journal of Information Sources and Services},
14(4):165--174, 2024.

\bibitem{bansal2022} Malti Bansal, Apoorva Goyal, and Apoorva Choudhary. Stock market prediction
with high accuracy using machine learning techniques. \textit{Procedia Computer Science},
215:247--265, 2022.

\bibitem{kompella2020} Subhadra Kompella and Kalyana Chakravarthy Chilukuri. Stock market
prediction using machine learning methods. \textit{International Journal of Computer Engineering
and Technology}, 10(3), 2019/2020.

\bibitem{ampomah2021} Ernest Kwame Ampomah, Gabriel Nyame, Zhiguang Qin, Prince Clement
Addo, Enoch Opanin Gyamfi, and Micheal Gyan. Stock market prediction with gaussian na\"ive bayes
machine learning algorithm. \textit{Informatica}, 45(2):243--256, 2021.

\end{thebibliography}

\end{document}
